\section{Особенности реализации}

\subsection{Общий поток выполнения}

Проект собирается как набор статических библиотек через CMake. Точка входа~---
единый исполняемый файл \texttt{graph\_main}, который инициализирует
\texttt{Menu} и запускает цикл ввода. Зависимости между модулями
однонаправленные: лабораторные зависят от \texttt{common}, \texttt{common}
не знает ни о каких лабораторных.

Типичный сценарий выполнения выглядит так. Пользователь выбирает пункт меню,
\texttt{Menu} вызывает соответствующий метод \texttt{Runner}'а из нужной лабы.
\texttt{Runner} читает параметры из консоли, делегирует работу алгоритмическому
классу (\texttt{Generator}, \texttt{ShimbellMethod} и т.\,д.), затем передаёт
результат в \texttt{FileHandler} для сохранения и в \texttt{Visualizer} для
отрисовки. \texttt{Visualizer} вызывает Python-скрипт через \texttt{system()},
передавая пути к файлам как аргументы командной строки. Python читает данные
из \texttt{assets/txt/} и сохраняет картинку в \texttt{assets/png/}.

\begin{figure}[h]
\centering
\begin{verbatim}
main.cc -> Menu -> Runner (lab1/lab2/...) -> Generator ->
    -> Algorithms -> FileHandler -> assets/txt/ ->
        -> Visualizer -> Python -> assets/png/
\end{verbatim}

\caption{Поток выполнения}
\end{figure}

\subsection{Модуль common}

Модуль содержит весь код, который переиспользуется несколькими лабораторными:
базовый граф, генератор, визуализатор, файловый обработчик, утилиты и меню.
Сборка через CMake в статическую библиотеку \texttt{graph\_common}.

\subsubsection{GraphBase: шаблонный граф}

\texttt{GraphBase<VertexT, EdgeT>}~--- базовый шаблонный класс, параметризованный
типами вершины и ребра.

\begin{lstlisting}[style=cppstyle, caption={Объявление GraphBase (GraphBase.h)}]
template <VertexType VertexT, EdgeType EdgeT>
class GraphBase {
    unordered_map<uint64_t, unique_ptr<VertexT>> m_vertices_;
    unordered_map<uint64_t, EdgeT> m_edges_;
    bool is_directed_ = false;
    ...
};
\end{lstlisting}

Класс служит основой для \texttt{Graph} (\texttt{GraphBase<Vertex, EdgeData>})
и \\ \texttt{FlowNetwork} (\texttt{GraphBase<FlowVertex, FlowEdge>}).

\paragraph{Хранение данных.} Вершины и рёбра лежат в хэш-таблицах с ключом
$\mathit{key}(from, to) = (from \ll 32)\;|\;to$. Для неориентированных
графов пара нормализуется перед хешированием: $from \leq to$. Это даёт
$O(1)$ проверку наличия ребра и автоматически обеспечивает симметрию.

\paragraph{Список инцидентности.} Смежные вершины хранятся в каждом объекте
вершины в \texttt{vector<int> m\_neighbors\_}. При вызове \texttt{addEdge}
обновляются и хэш-таблица рёбер, и список соседей. Запрос
\texttt{neighbors(v)} возвращает этот вектор за $O(\deg v)$. Такая
двойная структура~--- хэш по рёбрам плюс список соседей в вершине~---
позволяет не выбирать между скоростью обхода и скоростью проверки.

\paragraph{RAII и владение.} RAII (\textit{Resource Acquisition Is Initialization})~---
идиома C++, при которой ресурс захватывается в конструкторе и освобождается
в деструкторе, время жизни ресурса привязано к времени жизни объекта.
Вершины хранятся как \texttt{unique\_ptr<VertexT>}~---
владение однозначно, удаление происходит автоматически при разрушении
контейнера, утечки памяти исключены на уровне типа.

Ребро хранится как значение \texttt{EdgeT}~--- копирование дёшево
(структура из двух \texttt{int} и \texttt{double}). Вершина~--- указатель, потому что
\texttt{VertexT} содержит \texttt{vector<int> m\_neighbors\_},
который растёт динамически; хранение по значению в \texttt{unordered\_map}
привело бы к инвалидации ссылок при рехешировании.

Все немодифицирующие методы помечены \texttt{[[nodiscard]]}~---
атрибут стандарта C++17, при котором компилятор предупреждает об игнорировании
возвращаемого значения:

\begin{lstlisting}[style=cppstyle, caption={Публичный интерфейс GraphBase (GraphBase.h)}]
[[nodiscard]] bool hasVertex(int id) const noexcept;
[[nodiscard]] bool hasEdge(int from, int to)  const;
[[nodiscard]] std::optional<EdgeT> getEdge(int from, int to);
[[nodiscard]] std::vector<int> vertexIds() const;
[[nodiscard]] std::vector<EdgeT> edges() const;
[[nodiscard]] std::vector<int> neighbors(int id) const;
[[nodiscard]] int degree(int v) const;
\end{lstlisting}

Возвращаемые значения с неопределённым результатом (отсутствующее ребро)
возвращаются через \texttt{std::optional<EdgeT>}~--- явный способ
выразить «может не быть», без исключений и без неявного кода ошибки.

\subsubsection{Graph (DAG)}

\texttt{Graph} наследует \texttt{GraphBase<Vertex, EdgeData>}.

\begin{lstlisting}[style=cppstyle, caption={Graph, Vertex, EdgeData (Graph.h)}]
struct EdgeData {
    int from{-1};
    int to{-1};
    double weight{1.0};

    EdgeData() = default;
    EdgeData(int f, int t, double w) : from(f), to(t), weight(w) {}
    bool operator<(EdgeData const& other) const { return weight < other.weight; }
};

class Vertex {
public:
    explicit Vertex(int id) : m_id_(id) {}
    [[nodiscard]] int id() const noexcept { return m_id_; }
    void addNeighbor(int neighborId);
    [[nodiscard]] std::vector<int> neighbors() const;
private:
    int m_id_;
    std::vector<int> m_neighbors_;
};

class Graph : public GraphBase<Vertex, EdgeData> {
public:
    explicit Graph(bool isDirected = false) : GraphBase(isDirected) {}
    bool addEdge(int from, int to, double weight);
    [[nodiscard]] std::optional<double> getEdgeWeight(int from, int to) const;
    [[nodiscard]] std::vector<std::pair<int, double>> neighbors(int id) const;
    [[nodiscard]] std::vector<int> getNeighbors(int id) const;
    void printGraphInfo() const;
};
\end{lstlisting}

\texttt{EdgeData} хранит поля \texttt{from}, \texttt{to} и \texttt{weight}.
Оператор \texttt{<} по весу позволяет сортировать рёбра без дополнительных
компараторов~--- используется в алгоритмах, требующих упорядочения рёбер.

\texttt{Vertex} хранит идентификатор и список смежных вершин
\texttt{m\_neighbors\_}. Метод \texttt{addNeighbor} вызывается из
\texttt{Graph::addEdge}: для неориентированного графа смежность
добавляется в обоих направлениях~--- ребро регистрируется в хэш-таблице
через \texttt{GraphBase::addEdge}, а входящая вершина получает обратную
запись через явный вызов \texttt{m\_vertices\_[to]->addNeighbor(from)}.

\texttt{Graph} расширяет интерфейс базового класса тремя методами. \\
\texttt{getEdgeWeight(from, to)} возвращает \texttt{optional<double>}~---
вес ребра или \texttt{nullopt} если ребра нет.
\texttt{neighbors(id)} возвращает список смежных вершин с весами
\texttt{vector<pair<int,double>>}, объединяя данные из списка соседей
и хэш-таблицы рёбер за один проход.
\texttt{getNeighbors(id)} делегирует в  \\ \texttt{GraphBase::neighbors}
и возвращает только идентификаторы~--- используется там, где вес не нужен.

\subsubsection{Generator: генерация графов}

Генератор использует шаблонный приватный метод,
принимающий лямбды \texttt{addEdge} и \texttt{genWeight}. Это позволяет
переиспользовать один алгоритм структурной генерации для разных типов
графов и способов генерации весов. Публичные методы лишь собирают
нужные лямбды и передают их в шаблонный метод.

Для управления режимами используются \texttt{enum class}~--- исключают
неявное преобразование в \texttt{int}:

\begin{lstlisting}[style=cppstyle, caption={Перечисления параметров генератора (Generator.h)}]
enum class EdgeCountDist { Uniform, TruncatedNormal };
enum class WeightSign    { Positive, Negative, Mixed };
\end{lstlisting}

\paragraph{Степенная генерация: generateRiceGraphByDegrees.}
Метод строит граф в три шага, связанных через \texttt{buildByRice}.

\textbf{Шаг~1: число рёбер (\texttt{sampleEdgeCount}).}

\textbf{Вход:} $m_{\min} = n-1$, $m_{\max} = \frac{n(n-1)}{2}$,
режим \texttt{EdgeCountDist}.

\textbf{Выход:} \texttt{int}~--- число рёбер $m$.

\textbf{Описание:} \texttt{Uniform}~--- равномерная выборка из
$[m_{\min}, m_{\max}]$. \texttt{TruncatedNormal}~--- нормальное
распределение $\mathcal{N}(\mu, \sigma^2)$ с $\mu = (m_{\min} + m_{\max})/2$
и $\sigma = (m_{\max} - m_{\min})/6$, усечённое методом отклонения.

\begin{lstlisting}[style=cppstyle, caption={Выбор числа рёбер (Generator.cc)}]
int Generator::sampleEdgeCount(int minE,
                               int maxE,
                               EdgeCountDist dist)
{
    if (minE >= maxE) return minE;
    if (dist == EdgeCountDist::Uniform)
        return randomInt(minE, maxE);
    double mu    = (minE + maxE) / 2.0;
    double sigma = (maxE - minE) / 6.0;
    double val;

    do { val = mu + sigma * normal_dist_(rng_); }
    while (val < minE || val > maxE);

    return static_cast<int>(std::round(val));
}
\end{lstlisting}

\textbf{Шаг~2: целевые степени (\texttt{computeDegreesFromRice}).}

\textbf{Вход:} $n$, $a$, $h$, $m$.

\textbf{Выход:} \texttt{vector<int>} длины $n$.

\textbf{Описание:} Для каждой вершины сэмплируется
$x_i = \sqrt{(h + au_1)^2 + (au_2)^2}$, $u_1, u_2 \sim \mathcal{N}(0,1)$.
Вектор нормируется через softmax со стабилизацией. Степень вершины~$i$:
$q_i = p_i \cdot \mathit{cap}_i$, где $\mathit{cap}_i = n - 1 - i$~---
максимальная степень в топологическом порядке. Квоты масштабируются к
сумме $m$, целые части берутся через \texttt{floor}, остаток
распределяется жадно по убыванию дробных частей.

\begin{lstlisting}[style=cppstyle, caption={Вычисление целевых степеней (Generator.cc)}]
vector<int> Generator::computeDegreesFromRice(int n,
                                              bool isDirected,
                                              double a,
                                              double h,
                                              int totalEdges)
{
    std::vector<double> prob(n);
    for (auto& v : prob) v = riceWeight(a, h);

    double maxVal = *std::max_element(prob.begin(), prob.end());
    double sumExp = 0.0;
    for (auto& v : prob) { v = std::exp(v - maxVal); sumExp += v; }
    for (auto& v : prob) v /= sumExp;

    std::vector<int> cap(n);
    for (int i = 0; i < n; ++i) cap[i] = n - 1 - i;

    std::vector<double> quota(n);
    double sumW = 0.0;
    for (int i = 0; i < n; ++i) { quota[i] = prob[i] * cap[i]; sumW += quota[i]; }
    if (sumW > 0.0)
        for (auto& q : quota) q *= static_cast<double>(totalEdges) / sumW;

    std::vector<int> degrees(n);
    int allocated = 0;
    for (int i = 0; i < n; ++i) {
        degrees[i] = std::min((int)std::floor(quota[i]), cap[i]);
        allocated += degrees[i];
    }
    int remainder = totalEdges - allocated;
    std::vector<int> order(n);
    std::iota(order.begin(), order.end(), 0);
    std::sort(order.begin(), order.end(), [&](int a, int b) {
        return (quota[a] - std::floor(quota[a]))
             > (quota[b] - std::floor(quota[b]));
    });
    for (int idx : order) {
        if (remainder <= 0) break;
        if (degrees[idx] < cap[idx]) { ++degrees[idx]; --remainder; }
    }
    return degrees;
}
\end{lstlisting}

\textbf{Шаг~3: построение графа (\texttt{generateByDegreesTemplate}).}

\textbf{Вход:} $n$, вектор \texttt{degrees}, лямбды \texttt{addEdge}
и \texttt{genWeight}, флаг ориентированности.

\textbf{Выход:} \texttt{unique\_ptr<GraphT>}.

\textbf{Описание:} Первый проход строит остовное дерево: для вершины $i$
предок выбирается из вершин с ненулевым бюджетом степени; если таких нет~---
случайно. Второй проход добирает рёбра согласно оставшемуся бюджету.
Ациклические свойства гарантируются тем, что рёбра в обоих проходах
добавляются только от $u$ к $v$ при $u < v$.

\begin{lstlisting}[style=cppstyle, caption={Построение графа по целевым степеням (Generator.tpp)}]
template <typename GraphT,
          typename AddEdgeFunc,
          typename WeightGenFunc>
std::unique_ptr<GraphT> Generator::generateByDegreesTemplate(
                                    int numVertices,
                                    std::vector<int> degrees,
                                    AddEdgeFunc addEdge,
                                    WeightGenFunc genWeight,
                                    bool isDirected)
{
    auto graph = std::make_unique<GraphT>(isDirected);
    for (int i = 0; i < numVertices; ++i) graph->addVertex(i);

    for (int i = 1; i < numVertices; ++i) {
        std::vector<int> with_budget;
        for (int j = 0; j < i; ++j)
            if (degrees[j] > 0) with_budget.push_back(j);
        int from = with_budget.empty()
            ? randomInt(0, i - 1)
            : with_budget[randomInt(0, (int)with_budget.size() - 1)];
        addEdge(*graph, from, i, genWeight());
        if (degrees[from] > 0) --degrees[from];
    }

    for (int i = 0; i < numVertices - 1; ++i) {
        if (degrees[i] <= 0) continue;
        std::vector<int> available;
        for (int j = i + 1; j < numVertices; ++j)
            if (!graph->hasEdge(i, j)) available.push_back(j);
        std::shuffle(available.begin(), available.end(), rng_);
        int toAdd = std::min(degrees[i], (int)available.size());
        for (int k = 0; k < toAdd; ++k)
            addEdge(*graph, i, available[k], genWeight());
    }
    return graph;
}
\end{lstlisting}

Точка входа степенной генерации:

\begin{lstlisting}[style=cppstyle, caption={generateRiceGraphByDegrees (Generator.cc)}]
std::unique_ptr<Graph> Generator::generateRiceGraphByDegrees(
                                    int numVertices,
                                    bool isDirected,
                                    double a,
                                    double h,
                                    EdgeCountDist dist,
                                    WeightSign sign)
{
    auto gen_weight = [this, a, h, sign]() {
        return applySign(riceWeight(a, h), sign);
    };
    auto add_edge = [](Graph& g, int from, int to, double w) {
        g.addEdge(from, to, w);
    };
    return buildByRice<Graph>( numVertices, isDirected,
                               a, h, dist,
                               gen_weight, add_edge);
}
\end{lstlisting}

\subsubsection{CollectionUtils: шаблонные утилиты}

Статический класс с шаблонными методами для матриц и коллекций.

\texttt{makeMatrix<T>(rows, cols, getter)} строит матрицу любого типа,
заполняя её вызовами \texttt{getter(i, j)}. На его основе
\texttt{makeSquareMatrix} строит квадратную матрицу по списку вершин~---
используется в \texttt{FileHandler} для сохранения матриц смежности и весов.

\texttt{multiplyOptionalMatrix} выполняет матричное умножение над элементами типа \texttt{optional<double>}: \texttt{nullopt} означает отсутствие пути
и не участвует в агрегации. Компаратор передаётся параметром~---
одна функция работает как для $(\min,+)$, так и для $(\max,+)$
в методе Шимбелла.

\begin{lstlisting}[style=cppstyle, caption={Обобщённое матричное умножение (CollectionUtils.h)}]
template <typename CompareFunc>
static Matrix<std::optional<double>> multiplyOptionalMatrix(
    Matrix<std::optional<double>> const& a,
    Matrix<std::optional<double>> const& b,
    CompareFunc compare)
{
    int size = static_cast<int>(a.size());
    Matrix<optional<double>> result(
        size, vector<optional<double>>(size, nullopt)
    );
    for (int i = 0; i < size; ++i)
        for (int j = 0; j < size; ++j)
            for (int k = 0; k < size; ++k)
                if (a[i][k].has_value() && b[k][j].has_value()) {
                    double d = a[i][k].value() + b[k][j].value();
                    if (!result[i][j].has_value() ||
                        compare(d, result[i][j].value()))
                        result[i][j] = d;
                }
    return result;
}
\end{lstlisting}

\subsubsection{FileHandler: сохранение данных}

Статический класс. Все методы принимают имя файла и объект для сохранения.
\texttt{saveGraph} сохраняет список рёбер в формате \texttt{u v weight} по
одному ребру на строку~--- именно этот формат читает \texttt{GraphLoader}
в Python. \texttt{savePaths} пишет число путей первой строкой, затем
каждый путь как разделённую последовательность вершин.
Приватный шаблонный метод, предназначенный для чтение матрицы \texttt{saveSquareMatrix} используется публичными методами
\\ \texttt{saveAdjacencyMatrix}, \texttt{saveWeightMatrix},
\texttt{saveCapacityMatrix} и \texttt{saveCostMatrix}~--- все они
отличаются только геттером, который передаётся как лямбда.

\subsubsection{Visualizer: вызов Python}

Статический класс. Каждый метод формирует список строковых аргументов
и вызывает нужный скрипт через \texttt{runPythonScript}. Вызов идёт через \texttt{system()} с построенной
командной строкой.

\texttt{draw} вызывает \texttt{plot\_network.py} с флагом \texttt{--type graph}
или \texttt{--type flow}. \texttt{drawPaths}~--- \texttt{plot\_paths.py}
с дополнительным файлом путей. \texttt{drawMatrix}~--- \texttt{plot\_matrix.py}.
\texttt{drawColoredGraph}~--- \texttt{plot\_colored\_graph.py} с файлом
раскраски вершин.

\texttt{DrawDataConfig} хранит конфигурацию путей и заголовков для каждого
пункта меню в \texttt{map<int, DrawData>}. Все пути к файлам
собраны в одном месте, что упрощает добавление нового пункта меню.
